% YIN-VERBATIM: records=YIN-OY-T000016,YIN-OY-T000017; time=00:05:02.580-00:05:33.360
What is quantum field theory? This question is very easy to answer. A rough
answer would be: it's basically quantum mechanics, but with relativistic
locality. And \emph{[likely: I haven't]} explained what that means.

% EQUATION FROM NOTE PAGE 1 / PDF PAGE 6 HERE: Q: What is QFT?; A: QM + locality

% YIN-VERBATIM: records=YIN-OY-T000018,YIN-OY-T000019,YIN-OY-T000020; time=00:05:33.360-00:06:42.900
So, in fact, this is actually a confusing question even for the vast majority
of physicists, because I would say that there are really two conceptually
slightly different things that are both called quantum field theory, that are
really conceptually distinct. I want to discuss both in this course. One type
of quantum field theory is what I would call local, or sometimes referred to
as UV complete, and the other is what's known as effective field theory. So
what are these things? It's a little easier to state what a local QFT means.

% YIN-VERBATIM: records=YIN-OY-T000021; time=00:06:42.900-00:07:04.280
First of all, a local QFT is equipped with Poincar\'e symmetry. Poincar\'e
symmetry consists of translations in spacetime together with Lorentz
transformations.

% EQUATION FROM NOTE PAGE 1 / PDF PAGE 6 HERE: QM w/ Poincare sym; necessarily infinitely many d.o.f.

% YIN-VERBATIM: records=YIN-OY-T000022,YIN-OY-T000023,YIN-OY-T000024; time=00:07:04.280-00:08:04.800
Okay, so this kind of fundamental theory, first of all, is a quantum-mechanical
system. Some folks might have the slight misconception that quantum field
theory, going beyond quantum mechanics, is a generalization of quantum
mechanics. No, quantum field theory is not a generalization of quantum
mechanics. Quantum field theory is a specialization. Okay, so we are going to
allow quantum mechanics with infinitely many degrees of freedom, so infinitely
many\ldots{} The only requirement, well, not the only requirement, is one
requirement on top of the kind of quantum mechanics you learn in your
undergraduate quantum-mechanics class: we must also have Poincar\'e symmetry.
Now\ldots{}

% YIN-VERBATIM: records=YIN-OY-T000025,YIN-OY-T000026; time=00:08:04.800-00:08:31.620
\emph{[In a quantum]} mechanical system, a symmetry will be realized as a
unitary operator that acts on Hilbert space, and the operator will transform by
conjugation with the action of this unitary operator. So that will also be the
case here. I am giving a very broad overview for now. All of these words will
be put in the form of precise equations later.

% YIN-VERBATIM: records=YIN-OY-T000027; time=00:08:31.620-00:08:51.120
\textit{Yin:} All right. Any questions so far? You should feel free to
interrupt me at any point. You don't have to raise your hand in this class.

% YIN-VERBATIM: records=YIN-OY-T000028,YIN-OY-T000029,YIN-OY-T000030; time=00:08:51.120-00:09:38.279
But that's not all. We'll see that for a local QFT there also have to be local
operators, sometimes referred to as field operators. Well, we'll say what they
are. Importantly, these field operators have to obey some important
properties, including, particularly, microcausality. Now, basically, this is a
precise technical statement that captures the idea that signals cannot
propagate faster than the speed of light in a relativistic theory, but in a
precise quantum-mechanical setting. All right.

% EQUATION FROM NOTE PAGE 1 / PDF PAGE 6 HERE: local "field" operators; micro causality

% YIN-VERBATIM: records=YIN-OY-T000031,YIN-OY-T000032; time=00:09:38.279-00:10:02.160
Okay, now, what is effective field theory? Well, it's not really possible for
me to explain in detail, without setting up some technical tools, what
effective field theories are. At this point, I'll just say a few general facts
about it.

% YIN-VERBATIM: records=YIN-OY-T000033; time=00:10:02.160-00:10:10.980
It is typically\ldots{} To say \emph{typically}\ldots{} I like this. Thank
you. Okay, so\ldots{}

% YIN-VERBATIM: records=YIN-OY-T000034,YIN-OY-T000035,YIN-OY-T000036,YIN-OY-T000037,YIN-OY-T000038; time=00:10:10.980-00:11:38.880
Typically, \emph{[it is]} defined only perturbatively, in some coupling. So, in
contrast, \emph{[unclear]}, there are plenty of examples of quantum field
theories that are completely well-defined at strong coupling, and these
effective field theories typically are defined as a formal perturbative series
in the coupling constant, and it may not have \emph{[unclear]}. In particular,
it may or may not be convergent. And it's going to be based on a path integral,
a functional integral over the space of fields. So you can also use that to
define a local fundamental theory, but very often, depending on the details of
the theory, it may or may not make sense at the nonperturbative level. And
also\ldots{} An effective field theory typically, by design, captures
long-distance or low-energy observables, and not necessarily beyond that.

% EQUATION FROM NOTE PAGE 1 / PDF PAGE 6 HERE: defined perturbatively; path integral over space of fields; long-distance low-energy observables

% YIN-VERBATIM: records=YIN-OY-T000039; time=00:11:38.880-00:11:56.160
\textit{Yin:} So, what are some examples? Any questions about it?
\quad\textit{Student:} \emph{[unclear audio].}
\quad\textit{Yin:} Yes. I will describe that. Yeah.

% YIN-VERBATIM: records=YIN-OY-T000040,YIN-OY-T000041; time=00:11:56.160-00:12:22.380
So, let me just say very roughly, for a moment, that we'll have some field
operator. Let's say, let's call it \emph{[unclear; likely
$\hat\phi$]}, which will be associated with a spacetime point. Let's say it
depends on time coordinate $t$ and space coordinate $x$. Now\ldots{} To say
that you have an operator associated \emph{[to a point]}, by itself, is not a
very meaningful statement. I can parametrize \emph{[unclear; likely
$\hat\phi$]} however I want.

% YIN-VERBATIM: records=YIN-OY-T000042,YIN-OY-T000043,YIN-OY-T000044; time=00:12:22.380-00:13:17.820
Now, there are two key criteria, which I'll write as equations later, but I'll
just say them now, that say what it means to \emph{[be local]}. First of
all\ldots{} Under a Poincar\'e transformation, if you move a point to another
point in spacetime by a Poincar\'e-symmetry action, this corresponding unitary
operator acts on this \emph{[unclear; likely $\hat\phi$]} also by moving it to
\emph{[unclear; likely $\hat\phi$]} at the other point. Also, on top of that,
these operators have to obey microcausality, in the sense that you have
\emph{[unclear; likely $\hat\phi$]} at one point in spacetime and another at
another point in spacetime. These two points are spacelike separated and
causally disconnected; then the two operators commute. So that's kind of
microcausality. These two are the essential criteria for what it means for the
operator to be local. This is a concept that's rather alien if you only know
about undergraduate quantum mechanics, but that will be central to our study
of QFT.

% YIN-VERBATIM: records=YIN-OY-T000045; time=00:13:17.820-00:13:23.160
\textit{Yin:} Right. Any other questions?

% YIN-VERBATIM: records=YIN-OY-T000046; time=00:13:23.160-00:13:45.240
All right, so what are some examples? Let me write down some examples of local
QFT versus effective field theory, some of which we'll discuss in this course,
and \emph{[some]} not in this course; they will be in the next semester's
course.

% YIN-VERBATIM: records=YIN-OY-T000047,YIN-OY-T000048,YIN-OY-T000049,YIN-OY-T000050; time=00:13:45.240-00:14:40.680
So, there's the famous $\phi^4$ theory. This is the theory of a scalar field
with quartic interaction. We'll introduce this in a few weeks, actually. But
now, not in four-dimensional spacetime, but rather in $D$. This is the
spacetime dimension. Whenever I use capital $D$, I always mean the spacetime
dimension here, including time and space. So, if you construct such
a\ldots{} \ldots{}theory in two or three dimensions of Minkowski spacetime, so
either $1+1$D or $2+1$D, then that would be a completely fine local QFT. But on
the other hand, $\phi^4$ theory in four dimensions is only sensible, only makes
sense, at the level of effective field theory. All right.

% EQUATION FROM NOTE PAGE 1 / PDF PAGE 6 HERE: phi^4 theory in D=2,3; phi^4 theory in D=4

% YIN-VERBATIM: records=YIN-OY-T000051; time=00:14:40.680-00:15:00.000
Now, some of your favorite theories that you might want to study are,
unfortunately, actually not known to be well-defined at the nonperturbative
level as local QFTs\ldots{}

% YIN-VERBATIM-COUNT: eligible_cleaned_words=1036
% YIN-VERBATIM-COUNT: retained_spoken_words=1036
% YIN-VERBATIM-COUNT: retained_ratio=1.000000
